\section{Conclusão e Trabalhos Futuros}

Este trabalho apresentou o \textit{JEV-Cascade Reflection Loop}, uma arquitetura híbrida que integra LLMs geradores a um modelo discriminativo de decisão rápida (\textit{System One} -- JEV) para converter requisitos em linguagem natural irrestrita em templates semiformais MASTeR. Ao substituir a sequência de validadores autorregressivos (\textit{LLM-as-a-Judge}) e o filtro auxiliar de feedback por uma única consulta paralela tipada (\texttt{Choice}, \texttt{Noul} e \texttt{Score}) condicionada simultaneamente ao requisito fonte e ao candidato, a abordagem ataca diretamente o custo operacional e o fenômeno de \textit{over-editing} observados nos fluxos reflexivos recentes da literatura.

Além de eliminar por construção a necessidade de um LLM dedicado à filtragem de feedback não acionável, o uso de probabilidades calibradas viabiliza uma política explícita de governança humano-no-ciclo (\textit{Fast Accept}, \textit{Guided Rewrite} e \textit{Escalate}), essencial para especificações industriais de grande porte onde requisitos compostos permanecem desafiadores para automação plena. Como trabalhos futuros, pretende-se (i) calibrar dinamicamente os limiares $(\tau_{\text{rej}}, \tau_{\text{acc}})$ via predição conformal sobre avaliações de especialistas, (ii) estender o catálogo de primitivas JEV para outros templates como EARS e (iii) investigar a correlação direta do escore ordinal do JEV como métrica substituta de avaliação humana em escala.
